% ORIG03-SEGMENT-CODE: OR03-S002
% ORIG03-SEGMENT-ID: d90.orig.v1.tr03.seg0002
% ORIG03-MODULE-ID: diag.module.0001
% SPDX-License-Identifier: CC-BY-SA-4.0
% Karya asli berbahasa Indonesia; semua butir disusun independen.

\section{Diagnostik prasyarat}
\label{orig03:section:diagnostik}

Sepuluh butir berikut menguji prasyarat yang dipakai di seluruh buku. Setiap
subbutir merupakan prompt tersendiri dengan dua tingkat petunjuk, jawaban
ringkas, dan solusi lengkap. Materi ini tidak memuat pemodelan atau algoritma
yang menjadi wilayah O018.

% ORIG03-STABLE-ID: diag.problem.0001
\begin{exercise}[Norma dan ketaksamaan dasar]
\label{orig03:diag:problem:0001}
\begin{enumerate}
% ORIG03-STABLE-ID: diag.prompt.0001
\item\label{orig03:diag:prompt:0001}
Untuk $x=(3,-4)\in\mathbb{R}^2$, hitung $\lVert x\rVert_1$,
$\lVert x\rVert_2$, dan $\lVert x\rVert_\infty$.
% ORIG03-STABLE-ID: diag.prompt.0002
\item\label{orig03:diag:prompt:0002}
Buktikan bahwa untuk setiap $z\in\mathbb{R}^n$,
\[
  \lVert z\rVert_\infty\leq\lVert z\rVert_2
  \leq\lVert z\rVert_1\leq\sqrt n\,\lVert z\rVert_2.
\]
\end{enumerate}
\end{exercise}

% ORIG03-STABLE-ID: diag.hint1.0001
\par\noindent\textbf{Petunjuk tahap 1 (diag.prompt.0001).}
Gunakan langsung tiga definisi norma.
\label{orig03:diag:hint1:0001}
% ORIG03-STABLE-ID: diag.hint2.0001
\par\noindent\textbf{Petunjuk tahap 2 (diag.prompt.0001).}
Nilai mutlak koordinatnya adalah $3$ dan $4$.
\label{orig03:diag:hint2:0001}
% ORIG03-STABLE-ID: diag.answer.0001
\par\noindent\textbf{Jawaban ringkas (diag.prompt.0001).}
$\lVert x\rVert_1=7$, $\lVert x\rVert_2=5$, dan
$\lVert x\rVert_\infty=4$.
\label{orig03:diag:answer:0001}
% ORIG03-STABLE-ID: diag.solution.0001
\par\noindent\textbf{Solusi lengkap (diag.prompt.0001).}
Definisi memberi
$\lVert x\rVert_1=|3|+|-4|=7$,
$\lVert x\rVert_2=\sqrt{3^2+(-4)^2}=5$, dan
$\lVert x\rVert_\infty=\max\{3,4\}=4$.
\label{orig03:diag:solution:0001}

% ORIG03-STABLE-ID: diag.hint1.0002
\par\noindent\textbf{Petunjuk tahap 1 (diag.prompt.0002).}
Bandingkan setiap koordinat dengan jumlah kuadrat, lalu gunakan
Cauchy--Schwarz untuk ketaksamaan terakhir.
\label{orig03:diag:hint1:0002}
% ORIG03-STABLE-ID: diag.hint2.0002
\par\noindent\textbf{Petunjuk tahap 2 (diag.prompt.0002).}
Kembangkan $(\sum_i|z_i|)^2$ dan gunakan
$\sum_i|z_i|\leq\sqrt{\sum_i1^2}\sqrt{\sum_i z_i^2}$.
\label{orig03:diag:hint2:0002}
% ORIG03-STABLE-ID: diag.answer.0002
\par\noindent\textbf{Jawaban ringkas (diag.prompt.0002).}
Rantai tersebut berlaku untuk semua $z$; konstanta $\sqrt n$ tajam pada
vektor yang semua koordinat mutlaknya sama.
\label{orig03:diag:answer:0002}
% ORIG03-STABLE-ID: diag.solution.0002
\par\noindent\textbf{Solusi lengkap (diag.prompt.0002).}
Untuk setiap $i$, $|z_i|^2\leq\sum_j|z_j|^2$, sehingga
$\lVert z\rVert_\infty\leq\lVert z\rVert_2$. Selanjutnya,
$\sum_i|z_i|^2\leq(\sum_i|z_i|)^2$, sehingga
$\lVert z\rVert_2\leq\lVert z\rVert_1$. Akhirnya Cauchy--Schwarz memberi
$\sum_i|z_i|\leq(\sum_i1)^{1/2}(\sum_i|z_i|^2)^{1/2}$.
Ini menghasilkan ketaksamaan terakhir. Jika $|z_1|=\cdots=|z_n|>0$,
ketaksamaan terakhir menjadi kesamaan.
\emph{Rubrik bukti: \ref{orig03:rubric:proof:0001}.}
\label{orig03:diag:solution:0002}

% ORIG03-STABLE-ID: diag.problem.0002
\begin{exercise}[Proyeksi dan ortogonalitas]
\label{orig03:diag:problem:0002}
\begin{enumerate}
% ORIG03-STABLE-ID: diag.prompt.0003
\item\label{orig03:diag:prompt:0003}
Proyeksikan $y=(2,-1)$ pada hiperbidang
$H=\{x\in\mathbb{R}^2:x_1+x_2=0\}$.
% ORIG03-STABLE-ID: diag.prompt.0004
\item\label{orig03:diag:prompt:0004}
Untuk himpunan konveks tertutup takkosong $C$, buktikan bahwa
$p=P_C(y)$ jika dan hanya jika
$\langle y-p,z-p\rangle\leq0$ untuk setiap $z\in C$.
\end{enumerate}
\end{exercise}

% ORIG03-STABLE-ID: diag.hint1.0003
\par\noindent\textbf{Petunjuk tahap 1 (diag.prompt.0003).}
Gunakan normal $a=(1,1)$ dan rumus proyeksi pada $a^Tx=0$.
\label{orig03:diag:hint1:0003}
% ORIG03-STABLE-ID: diag.hint2.0003
\par\noindent\textbf{Petunjuk tahap 2 (diag.prompt.0003).}
Kurangkan $((a^Ty)/(a^Ta))a$ dari $y$.
\label{orig03:diag:hint2:0003}
% ORIG03-STABLE-ID: diag.answer.0003
\par\noindent\textbf{Jawaban ringkas (diag.prompt.0003).}
$P_H(y)=(3/2,-3/2)$.
\label{orig03:diag:answer:0003}
% ORIG03-STABLE-ID: diag.solution.0003
\par\noindent\textbf{Solusi lengkap (diag.prompt.0003).}
Karena $a^Ty=1$ dan $a^Ta=2$,
\[
P_H(y)=y-\frac{a^Ty}{a^Ta}a=(2,-1)-\frac12(1,1)
=\left(\frac32,-\frac32\right).
\]
Jumlah koordinat hasilnya nol, sedangkan selisih $y-P_H(y)$ sejajar normal
$a$, jadi hasil tersebut memang proyeksi ortogonal.
\label{orig03:diag:solution:0003}

% ORIG03-STABLE-ID: diag.hint1.0004
\par\noindent\textbf{Petunjuk tahap 1 (diag.prompt.0004).}
Untuk arah maju gunakan titik $p+t(z-p)$; untuk arah balik kembangkan
$\lVert y-z\rVert^2$.
\label{orig03:diag:hint1:0004}
% ORIG03-STABLE-ID: diag.hint2.0004
\par\noindent\textbf{Petunjuk tahap 2 (diag.prompt.0004).}
Minimalitas fungsi kuadrat dalam $t\in[0,1]$ memberi turunan kanan taknegatif
di $t=0$.
\label{orig03:diag:hint2:0004}
% ORIG03-STABLE-ID: diag.answer.0004
\par\noindent\textbf{Jawaban ringkas (diag.prompt.0004).}
Kondisi hasil kali dalam tersebut tepat merupakan kondisi optimalitas orde
pertama bagi minimisasi jarak kuadrat pada $C$.
\label{orig03:diag:answer:0004}
% ORIG03-STABLE-ID: diag.solution.0004
\par\noindent\textbf{Solusi lengkap (diag.prompt.0004).}
Jika $p=P_C(y)$, maka untuk $z\in C$ dan $t\in[0,1]$, kekonveksan memberi
$p+t(z-p)\in C$. Fungsi
$q(t)=\lVert y-p-t(z-p)\rVert^2$ minimum di $t=0$, sehingga
$q'(0)=-2\langle y-p,z-p\rangle\geq0$. Sebaliknya, bila kondisi hasil kali
dalam berlaku, maka untuk setiap $z\in C$,
\[
\lVert y-z\rVert^2
=\lVert y-p\rVert^2+\lVert z-p\rVert^2
-2\langle y-p,z-p\rangle\geq\lVert y-p\rVert^2.
\]
Jadi $p$ meminimalkan jarak. Keunikan mengikuti kekonveksan ketat jarak
kuadrat.
\emph{Rubrik bukti: \ref{orig03:rubric:proof:0004}.}
\label{orig03:diag:solution:0004}

% ORIG03-STABLE-ID: diag.problem.0003
\begin{exercise}[Kombinasi dan irisan konveks]
\label{orig03:diag:problem:0003}
\begin{enumerate}
% ORIG03-STABLE-ID: diag.prompt.0005
\item\label{orig03:diag:prompt:0005}
Tentukan $t\in[0,1]$ sehingga
$(1,2)=t(0,1)+(1-t)(2,3)$.
% ORIG03-STABLE-ID: diag.prompt.0006
\item\label{orig03:diag:prompt:0006}
Buktikan bahwa irisan sebarang keluarga himpunan konveks adalah konveks.
\end{enumerate}
\end{exercise}

% ORIG03-STABLE-ID: diag.hint1.0005
\par\noindent\textbf{Petunjuk tahap 1 (diag.prompt.0005).}
Samakan koordinat pertama.
\label{orig03:diag:hint1:0005}
% ORIG03-STABLE-ID: diag.hint2.0005
\par\noindent\textbf{Petunjuk tahap 2 (diag.prompt.0005).}
Persamaan $2(1-t)=1$ memberi kandidat tunggal; periksa koordinat kedua.
\label{orig03:diag:hint2:0005}
% ORIG03-STABLE-ID: diag.answer.0005
\par\noindent\textbf{Jawaban ringkas (diag.prompt.0005).}
$t=1/2$.
\label{orig03:diag:answer:0005}
% ORIG03-STABLE-ID: diag.solution.0005
\par\noindent\textbf{Solusi lengkap (diag.prompt.0005).}
Koordinat pertama memberi $2(1-t)=1$, jadi $t=1/2$. Koordinat kedua lalu
bernilai $t+3(1-t)=1/2+3/2=2$. Karena $1/2\in[0,1]$, ini kombinasi konveks
yang sah.
\label{orig03:diag:solution:0005}

% ORIG03-STABLE-ID: diag.hint1.0006
\par\noindent\textbf{Petunjuk tahap 1 (diag.prompt.0006).}
Ambil dua titik dalam irisan dan uji kombinasi konveksnya pada setiap anggota
keluarga.
\label{orig03:diag:hint1:0006}
% ORIG03-STABLE-ID: diag.hint2.0006
\par\noindent\textbf{Petunjuk tahap 2 (diag.prompt.0006).}
Keanggotaan pada irisan berarti keanggotaan serentak pada semua himpunan.
\label{orig03:diag:hint2:0006}
% ORIG03-STABLE-ID: diag.answer.0006
\par\noindent\textbf{Jawaban ringkas (diag.prompt.0006).}
Irisan tetap konveks; irisan kosong juga konveks secara vakum.
\label{orig03:diag:answer:0006}
% ORIG03-STABLE-ID: diag.solution.0006
\par\noindent\textbf{Solusi lengkap (diag.prompt.0006).}
Tuliskan $C=\bigcap_{i\in I}C_i$ dengan setiap $C_i$ konveks. Jika
$x,y\in C$ dan $t\in[0,1]$, maka $x,y\in C_i$ untuk setiap $i$. Karena
$C_i$ konveks, $tx+(1-t)y\in C_i$ untuk setiap $i$, sehingga titik itu berada
di $C$. Bila $C$ kosong, implikasi dalam definisi konveks benar secara vakum.
\emph{Rubrik bukti: \ref{orig03:rubric:proof:0001}.}
\label{orig03:diag:solution:0006}

% ORIG03-STABLE-ID: diag.problem.0004
\begin{exercise}[Kuadratik, gradien, dan spektrum]
\label{orig03:diag:problem:0004}
\begin{enumerate}
% ORIG03-STABLE-ID: diag.prompt.0007
\item\label{orig03:diag:prompt:0007}
Untuk $f(x)=\tfrac12x^TQx+b^Tx$ dengan
$Q=\begin{psmallmatrix}4&1\\1&2\end{psmallmatrix}$ dan $b=(-1,3)$,
hitung $\nabla f(1,-1)$ dan $\nabla^2f$.
% ORIG03-STABLE-ID: diag.prompt.0008
\item\label{orig03:diag:prompt:0008}
Tentukan konstanta kuat-konveks terbesar $\mu$ dan konstanta Lipschitz gradien
terkecil $L$ bagi $f$.
\end{enumerate}
\end{exercise}

% ORIG03-STABLE-ID: diag.hint1.0007
\par\noindent\textbf{Petunjuk tahap 1 (diag.prompt.0007).}
Karena $Q$ simetris, $\nabla f(x)=Qx+b$.
\label{orig03:diag:hint1:0007}
% ORIG03-STABLE-ID: diag.hint2.0007
\par\noindent\textbf{Petunjuk tahap 2 (diag.prompt.0007).}
Hitung $Q(1,-1)^T=(3,-1)^T$.
\label{orig03:diag:hint2:0007}
% ORIG03-STABLE-ID: diag.answer.0007
\par\noindent\textbf{Jawaban ringkas (diag.prompt.0007).}
$\nabla f(1,-1)=(2,2)$ dan $\nabla^2f=Q$.
\label{orig03:diag:answer:0007}
% ORIG03-STABLE-ID: diag.solution.0007
\par\noindent\textbf{Solusi lengkap (diag.prompt.0007).}
Diferensiasi bentuk kuadratik simetris memberi $\nabla f(x)=Qx+b$ dan
$\nabla^2f=Q$. Pada $x=(1,-1)$,
$Qx=(4-1,1-2)=(3,-1)$, sehingga $Qx+b=(2,2)$.
\label{orig03:diag:solution:0007}

% ORIG03-STABLE-ID: diag.hint1.0008
\par\noindent\textbf{Petunjuk tahap 1 (diag.prompt.0008).}
Cari kedua nilai eigen matriks simetris $Q$.
\label{orig03:diag:hint1:0008}
% ORIG03-STABLE-ID: diag.hint2.0008
\par\noindent\textbf{Petunjuk tahap 2 (diag.prompt.0008).}
Persamaan karakteristiknya $\lambda^2-6\lambda+7=0$.
\label{orig03:diag:hint2:0008}
% ORIG03-STABLE-ID: diag.answer.0008
\par\noindent\textbf{Jawaban ringkas (diag.prompt.0008).}
$\mu=3-\sqrt2$ dan $L=3+\sqrt2$.
\label{orig03:diag:answer:0008}
% ORIG03-STABLE-ID: diag.solution.0008
\par\noindent\textbf{Solusi lengkap (diag.prompt.0008).}
Nilai eigen $Q$ adalah akar
$\lambda^2-6\lambda+7=0$, yaitu $3\pm\sqrt2$. Untuk kuadratik dengan Hessian
konstan simetris positif definit, konstanta kuat-konveks terbesar adalah nilai
eigen minimum dan konstanta Lipschitz gradien terkecil adalah nilai eigen
maksimum. Jadi hasilnya $\mu=3-\sqrt2$ dan $L=3+\sqrt2$.
\emph{Rubrik bukti: \ref{orig03:rubric:proof:0002}.}
\label{orig03:diag:solution:0008}

% ORIG03-STABLE-ID: diag.problem.0005
\begin{exercise}[Kondisi orde pertama]
\label{orig03:diag:problem:0005}
\begin{enumerate}
% ORIG03-STABLE-ID: diag.prompt.0009
\item\label{orig03:diag:prompt:0009}
Untuk $f(x)=(x-2)^2+1$, tuliskan penyangga afin orde pertama di $x_0=0$
dan verifikasi bahwa penyangga itu menyentuh $f$ di $x_0$.
% ORIG03-STABLE-ID: diag.prompt.0010
\item\label{orig03:diag:prompt:0010}
Buktikan bahwa jika $f$ konveks dan terdiferensialkan serta
$\nabla f(x^*)=0$, maka $x^*$ adalah minimizer global.
\end{enumerate}
\end{exercise}

% ORIG03-STABLE-ID: diag.hint1.0009
\par\noindent\textbf{Petunjuk tahap 1 (diag.prompt.0009).}
Gunakan $f(x_0)+f'(x_0)(x-x_0)$.
\label{orig03:diag:hint1:0009}
% ORIG03-STABLE-ID: diag.hint2.0009
\par\noindent\textbf{Petunjuk tahap 2 (diag.prompt.0009).}
$f(0)=5$ dan $f'(0)=-4$.
\label{orig03:diag:hint2:0009}
% ORIG03-STABLE-ID: diag.answer.0009
\par\noindent\textbf{Jawaban ringkas (diag.prompt.0009).}
Penyangga afinnya $\ell(x)=5-4x$, dan $\ell(0)=f(0)=5$.
\label{orig03:diag:answer:0009}
% ORIG03-STABLE-ID: diag.solution.0009
\par\noindent\textbf{Solusi lengkap (diag.prompt.0009).}
Turunan $f'(x)=2(x-2)$, jadi $f'(0)=-4$. Karena itu
$\ell(x)=f(0)+f'(0)x=5-4x$. Selisihnya
$f(x)-\ell(x)=x^2\geq0$, dan selisih nol di $x=0$.
\label{orig03:diag:solution:0009}

% ORIG03-STABLE-ID: diag.hint1.0010
\par\noindent\textbf{Petunjuk tahap 1 (diag.prompt.0010).}
Gunakan ketaksamaan orde pertama bagi fungsi konveks.
\label{orig03:diag:hint1:0010}
% ORIG03-STABLE-ID: diag.hint2.0010
\par\noindent\textbf{Petunjuk tahap 2 (diag.prompt.0010).}
Substitusikan gradien nol dalam
$f(y)\geq f(x^*)+\langle\nabla f(x^*),y-x^*\rangle$.
\label{orig03:diag:hint2:0010}
% ORIG03-STABLE-ID: diag.answer.0010
\par\noindent\textbf{Jawaban ringkas (diag.prompt.0010).}
$f(y)\geq f(x^*)$ untuk setiap $y$.
\label{orig03:diag:answer:0010}
% ORIG03-STABLE-ID: diag.solution.0010
\par\noindent\textbf{Solusi lengkap (diag.prompt.0010).}
Kekonveksan dan keterdiferensialan memberi, untuk setiap $y$,
\[
f(y)\geq f(x^*)+\langle\nabla f(x^*),y-x^*\rangle=f(x^*).
\]
Jadi tidak ada titik dengan nilai lebih kecil daripada $x^*$.
\emph{Rubrik bukti: \ref{orig03:rubric:proof:0001}.}
\label{orig03:diag:solution:0010}

% ORIG03-STABLE-ID: diag.problem.0006
\begin{exercise}[Subgradien dan operator proksimal]
\label{orig03:diag:problem:0006}
\begin{enumerate}
% ORIG03-STABLE-ID: diag.prompt.0011
\item\label{orig03:diag:prompt:0011}
Tentukan $\partial|x|$ pada $x=-2$, $x=0$, dan $x=3$.
% ORIG03-STABLE-ID: diag.prompt.0012
\item\label{orig03:diag:prompt:0012}
Hitung $\operatorname{prox}_{|\cdot|}(-3)$.
\end{enumerate}
\end{exercise}

% ORIG03-STABLE-ID: diag.hint1.0011
\par\noindent\textbf{Petunjuk tahap 1 (diag.prompt.0011).}
Di luar nol gunakan turunan; di nol gunakan definisi subgradien.
\label{orig03:diag:hint1:0011}
% ORIG03-STABLE-ID: diag.hint2.0011
\par\noindent\textbf{Petunjuk tahap 2 (diag.prompt.0011).}
Pada nol, $g$ harus memenuhi $|y|\geq gy$ untuk semua $y$.
\label{orig03:diag:hint2:0011}
% ORIG03-STABLE-ID: diag.answer.0011
\par\noindent\textbf{Jawaban ringkas (diag.prompt.0011).}
$\partial|-2|=\{-1\}$, $\partial|0|=[-1,1]$, dan
$\partial|3|=\{1\}$.
\label{orig03:diag:answer:0011}
% ORIG03-STABLE-ID: diag.solution.0011
\par\noindent\textbf{Solusi lengkap (diag.prompt.0011).}
Fungsi $|x|$ terdiferensialkan dengan turunan $-1$ untuk $x<0$ dan $1$
untuk $x>0$. Pada nol, syarat subgradien adalah $|y|\geq gy$ untuk setiap
$y$. Pemilihan $y>0$ memberi $g\leq1$, sedangkan $y<0$ memberi $g\geq-1$.
Jadi subdiferensial di nol adalah $[-1,1]$.
\label{orig03:diag:solution:0011}

% ORIG03-STABLE-ID: diag.hint1.0012
\par\noindent\textbf{Petunjuk tahap 1 (diag.prompt.0012).}
Gunakan rumus ambang lunak dengan ambang satu.
\label{orig03:diag:hint1:0012}
% ORIG03-STABLE-ID: diag.hint2.0012
\par\noindent\textbf{Petunjuk tahap 2 (diag.prompt.0012).}
$S_1(v)=\operatorname{sign}(v)\max\{|v|-1,0\}$.
\label{orig03:diag:hint2:0012}
% ORIG03-STABLE-ID: diag.answer.0012
\par\noindent\textbf{Jawaban ringkas (diag.prompt.0012).}
$\operatorname{prox}_{|\cdot|}(-3)=-2$.
\label{orig03:diag:answer:0012}
% ORIG03-STABLE-ID: diag.solution.0012
\par\noindent\textbf{Solusi lengkap (diag.prompt.0012).}
Operator proksimal meminimalkan $|x|+\tfrac12(x+3)^2$. Karena minimizer
negatif, turunan pada cabang $x<0$ adalah $-1+x+3=x+2$, yang nol pada
$x=-2$. Titik itu konsisten dengan cabangnya dan memberi hasil yang sama
dengan $S_1(-3)=-2$.
\label{orig03:diag:solution:0012}

% ORIG03-STABLE-ID: diag.problem.0007
\begin{exercise}[Ekspektasi dan rata-rata sampel]
\label{orig03:diag:problem:0007}
\begin{enumerate}
% ORIG03-STABLE-ID: diag.prompt.0013
\item\label{orig03:diag:prompt:0013}
Peubah acak $X$ bernilai $-1$ dengan peluang $2/3$ dan $2$ dengan peluang
$1/3$. Hitung $\mathbb{E}X$ dan $\operatorname{Var}(X)$.
% ORIG03-STABLE-ID: diag.prompt.0014
\item\label{orig03:diag:prompt:0014}
Jika $X_1,\dots,X_m$ saling independen dan berdistribusi seperti $X$, tentukan
ekspektasi dan varians $\bar X_m=m^{-1}\sum_{i=1}^mX_i$.
\end{enumerate}
\end{exercise}

% ORIG03-STABLE-ID: diag.hint1.0013
\par\noindent\textbf{Petunjuk tahap 1 (diag.prompt.0013).}
Hitung momen pertama dan kedua.
\label{orig03:diag:hint1:0013}
% ORIG03-STABLE-ID: diag.hint2.0013
\par\noindent\textbf{Petunjuk tahap 2 (diag.prompt.0013).}
$\mathbb{E}X^2=(2/3)1+(1/3)4$.
\label{orig03:diag:hint2:0013}
% ORIG03-STABLE-ID: diag.answer.0013
\par\noindent\textbf{Jawaban ringkas (diag.prompt.0013).}
$\mathbb{E}X=0$ dan $\operatorname{Var}(X)=2$.
\label{orig03:diag:answer:0013}
% ORIG03-STABLE-ID: diag.solution.0013
\par\noindent\textbf{Solusi lengkap (diag.prompt.0013).}
Kita memperoleh
$\mathbb{E}X=(2/3)(-1)+(1/3)2=0$ dan
$\mathbb{E}X^2=(2/3)1+(1/3)4=2$. Maka
$\operatorname{Var}(X)=\mathbb{E}X^2-(\mathbb{E}X)^2=2$.
\label{orig03:diag:solution:0013}

% ORIG03-STABLE-ID: diag.hint1.0014
\par\noindent\textbf{Petunjuk tahap 1 (diag.prompt.0014).}
Gunakan linearitas ekspektasi dan aditivitas varians bagi peubah independen.
\label{orig03:diag:hint1:0014}
% ORIG03-STABLE-ID: diag.hint2.0014
\par\noindent\textbf{Petunjuk tahap 2 (diag.prompt.0014).}
Faktor $1/m$ dikuadratkan ketika menghitung varians.
\label{orig03:diag:hint2:0014}
% ORIG03-STABLE-ID: diag.answer.0014
\par\noindent\textbf{Jawaban ringkas (diag.prompt.0014).}
$\mathbb{E}\bar X_m=0$ dan $\operatorname{Var}(\bar X_m)=2/m$.
\label{orig03:diag:answer:0014}
% ORIG03-STABLE-ID: diag.solution.0014
\par\noindent\textbf{Solusi lengkap (diag.prompt.0014).}
Linearitas memberi
$\mathbb{E}\bar X_m=m^{-1}\sum_i\mathbb{E}X_i=0$. Independensi memberi
\[
\operatorname{Var}(\bar X_m)=\frac1{m^2}\sum_{i=1}^m
\operatorname{Var}(X_i)=\frac1{m^2}(2m)=\frac2m.
\]
\emph{Rubrik bukti: \ref{orig03:rubric:proof:0003}.}
\label{orig03:diag:solution:0014}

% ORIG03-STABLE-ID: diag.problem.0008
\begin{exercise}[Rekurensi kontraktif]
\label{orig03:diag:problem:0008}
\begin{enumerate}
% ORIG03-STABLE-ID: diag.prompt.0015
\item\label{orig03:diag:prompt:0015}
Selesaikan $x_{k+1}=\tfrac34x_k+1$ dengan $x_0=0$ dan tentukan titik
tetapnya.
% ORIG03-STABLE-ID: diag.prompt.0016
\item\label{orig03:diag:prompt:0016}
Tentukan bilangan bulat terkecil $k$ yang menjamin $|x_k-4|\leq10^{-2}$.
\end{enumerate}
\end{exercise}

% ORIG03-STABLE-ID: diag.hint1.0015
\par\noindent\textbf{Petunjuk tahap 1 (diag.prompt.0015).}
Kurangkan titik tetap dari kedua ruas.
\label{orig03:diag:hint1:0015}
% ORIG03-STABLE-ID: diag.hint2.0015
\par\noindent\textbf{Petunjuk tahap 2 (diag.prompt.0015).}
Titik tetap memenuhi $x=\tfrac34x+1$.
\label{orig03:diag:hint2:0015}
% ORIG03-STABLE-ID: diag.answer.0015
\par\noindent\textbf{Jawaban ringkas (diag.prompt.0015).}
Titik tetapnya $4$ dan $x_k=4(1-(3/4)^k)$.
\label{orig03:diag:answer:0015}
% ORIG03-STABLE-ID: diag.solution.0015
\par\noindent\textbf{Solusi lengkap (diag.prompt.0015).}
Persamaan titik tetap memberi $x^*=4$. Dengan $e_k=x_k-4$ diperoleh
$e_{k+1}=\tfrac34e_k$ dan $e_0=-4$. Jadi
$e_k=-4(3/4)^k$ dan $x_k=4(1-(3/4)^k)$.
\label{orig03:diag:solution:0015}

% ORIG03-STABLE-ID: diag.hint1.0016
\par\noindent\textbf{Petunjuk tahap 1 (diag.prompt.0016).}
Gunakan bentuk tertutup galat dari subbutir sebelumnya.
\label{orig03:diag:hint1:0016}
% ORIG03-STABLE-ID: diag.hint2.0016
\par\noindent\textbf{Petunjuk tahap 2 (diag.prompt.0016).}
Selesaikan $(3/4)^k\leq1/400$ dengan logaritma.
\label{orig03:diag:hint2:0016}
% ORIG03-STABLE-ID: diag.answer.0016
\par\noindent\textbf{Jawaban ringkas (diag.prompt.0016).}
$k=21$.
\label{orig03:diag:answer:0016}
% ORIG03-STABLE-ID: diag.solution.0016
\par\noindent\textbf{Solusi lengkap (diag.prompt.0016).}
Karena $|x_k-4|=4(3/4)^k$, syaratnya adalah
$(3/4)^k\leq1/400$. Karena $\log(3/4)<0$,
\[
k\geq\frac{\log(1/400)}{\log(3/4)}\approx20.827.
\]
Bilangan bulat terkecil yang memenuhi adalah $21$.
\label{orig03:diag:solution:0016}

% ORIG03-STABLE-ID: diag.problem.0009
\begin{exercise}[Kerucut normal dan titik tetap proyeksi]
\label{orig03:diag:problem:0009}
\begin{enumerate}
% ORIG03-STABLE-ID: diag.prompt.0017
\item\label{orig03:diag:prompt:0017}
Minimalkan $f(x)=(x-2)^2$ pada $C=[0,1]$ dan verifikasi
$0\in f'(x^*)+N_C(x^*)$ dengan konvensi
$N_C(x)=\{v:v(y-x)\leq0\ \forall y\in C\}$.
% ORIG03-STABLE-ID: diag.prompt.0018
\item\label{orig03:diag:prompt:0018}
Verifikasi bahwa untuk setiap $\gamma>0$, solusi tersebut memenuhi
$x^*=P_C(x^*-\gamma f'(x^*))$.
\end{enumerate}
\end{exercise}

% ORIG03-STABLE-ID: diag.hint1.0017
\par\noindent\textbf{Petunjuk tahap 1 (diag.prompt.0017).}
Minimizer tak terkendala berada di luar interval.
\label{orig03:diag:hint1:0017}
% ORIG03-STABLE-ID: diag.hint2.0017
\par\noindent\textbf{Petunjuk tahap 2 (diag.prompt.0017).}
Pada ujung kanan, $N_{[0,1]}(1)=[0,\infty)$.
\label{orig03:diag:hint2:0017}
% ORIG03-STABLE-ID: diag.answer.0017
\par\noindent\textbf{Jawaban ringkas (diag.prompt.0017).}
$x^*=1$, $f'(1)=-2$, dan $2\in N_C(1)$.
\label{orig03:diag:answer:0017}
% ORIG03-STABLE-ID: diag.solution.0017
\par\noindent\textbf{Solusi lengkap (diag.prompt.0017).}
Fungsi menurun pada $[0,1]$, jadi minimumnya tercapai di $x^*=1$. Turunannya
$f'(x)=2(x-2)$ memberi $f'(1)=-2$. Dari definisi, setiap $v\geq0$ memenuhi
$v(y-1)\leq0$ untuk $y\in[0,1]$, sehingga $N_C(1)=[0,\infty)$. Memilih
$v=2$ menghasilkan $0=-2+2\in f'(1)+N_C(1)$.
\emph{Rubrik bukti: \ref{orig03:rubric:proof:0004}.}
\label{orig03:diag:solution:0017}

% ORIG03-STABLE-ID: diag.hint1.0018
\par\noindent\textbf{Petunjuk tahap 1 (diag.prompt.0018).}
Substitusikan $x^*=1$ dan $f'(1)=-2$.
\label{orig03:diag:hint1:0018}
% ORIG03-STABLE-ID: diag.hint2.0018
\par\noindent\textbf{Petunjuk tahap 2 (diag.prompt.0018).}
Argumen proyeksinya menjadi $1+2\gamma>1$.
\label{orig03:diag:hint2:0018}
% ORIG03-STABLE-ID: diag.answer.0018
\par\noindent\textbf{Jawaban ringkas (diag.prompt.0018).}
$P_{[0,1]}(1+2\gamma)=1=x^*$.
\label{orig03:diag:answer:0018}
% ORIG03-STABLE-ID: diag.solution.0018
\par\noindent\textbf{Solusi lengkap (diag.prompt.0018).}
Untuk $\gamma>0$,
$x^*-\gamma f'(x^*)=1-\gamma(-2)=1+2\gamma$. Proyeksi setiap bilangan lebih
besar daripada satu pada $[0,1]$ adalah satu. Karena itu
$P_C(x^*-\gamma f'(x^*))=1=x^*$.
\label{orig03:diag:solution:0018}

% ORIG03-STABLE-ID: diag.problem.0010
\begin{exercise}[Cauchy--Schwarz dan Young]
\label{orig03:diag:problem:0010}
\begin{enumerate}
% ORIG03-STABLE-ID: diag.prompt.0019
\item\label{orig03:diag:prompt:0019}
Untuk $u=(1,2)$ dan $v=(2,-1)$, hitung hasil kali dalam dan kedua normanya,
lalu periksa ketaksamaan Cauchy--Schwarz.
% ORIG03-STABLE-ID: diag.prompt.0020
\item\label{orig03:diag:prompt:0020}
Buktikan bahwa untuk $a,b\in\mathbb{R}$ dan $\eta>0$,
$2ab\leq\eta a^2+\eta^{-1}b^2$, serta tentukan syarat kesamaan.
\end{enumerate}
\end{exercise}

% ORIG03-STABLE-ID: diag.hint1.0019
\par\noindent\textbf{Petunjuk tahap 1 (diag.prompt.0019).}
Hitung $u^Tv$ sebelum normanya.
\label{orig03:diag:hint1:0019}
% ORIG03-STABLE-ID: diag.hint2.0019
\par\noindent\textbf{Petunjuk tahap 2 (diag.prompt.0019).}
Kedua norma sama dengan $\sqrt5$.
\label{orig03:diag:hint2:0019}
% ORIG03-STABLE-ID: diag.answer.0019
\par\noindent\textbf{Jawaban ringkas (diag.prompt.0019).}
$\langle u,v\rangle=0$ dan $0\leq5$.
\label{orig03:diag:answer:0019}
% ORIG03-STABLE-ID: diag.solution.0019
\par\noindent\textbf{Solusi lengkap (diag.prompt.0019).}
$\langle u,v\rangle=1\cdot2+2\cdot(-1)=0$, sedangkan
$\lVert u\rVert_2=\lVert v\rVert_2=\sqrt5$. Jadi
$|\langle u,v\rangle|=0\leq\lVert u\rVert_2\lVert v\rVert_2=5$.
Vektor-vektor tersebut ortogonal, bukan sejajar, sehingga ketaksamaan tidak
menjadi kesamaan pada batas atas yang positif.
\label{orig03:diag:solution:0019}

% ORIG03-STABLE-ID: diag.hint1.0020
\par\noindent\textbf{Petunjuk tahap 1 (diag.prompt.0020).}
Mulai dari kuadrat yang taknegatif.
\label{orig03:diag:hint1:0020}
% ORIG03-STABLE-ID: diag.hint2.0020
\par\noindent\textbf{Petunjuk tahap 2 (diag.prompt.0020).}
Kembangkan $(\sqrt\eta\,a-b/\sqrt\eta)^2$.
\label{orig03:diag:hint2:0020}
% ORIG03-STABLE-ID: diag.answer.0020
\par\noindent\textbf{Jawaban ringkas (diag.prompt.0020).}
Ketaksamaan selalu berlaku dan kesamaan terjadi tepat ketika $b=\eta a$.
\label{orig03:diag:answer:0020}
% ORIG03-STABLE-ID: diag.solution.0020
\par\noindent\textbf{Solusi lengkap (diag.prompt.0020).}
Karena
\[
0\leq\left(\sqrt\eta\,a-\frac{b}{\sqrt\eta}\right)^2
=\eta a^2-2ab+\eta^{-1}b^2,
\]
pemindahan suku memberi ketaksamaan yang diminta. Kuadrat itu nol tepat jika
$\sqrt\eta\,a=b/\sqrt\eta$, yaitu $b=\eta a$.
\emph{Rubrik bukti: \ref{orig03:rubric:proof:0001}.}
\label{orig03:diag:solution:0020}
